\subsubsection{同性伴侣的法律事务：现行法框架下的可用工具}

婚姻登记是异性伴侣专有的制度安排，但这不意味着同性伴侣在财产、医疗与身后事务上完全无工具可用。本节只做一件事：\textbf{把现行法律框架下实际可用的工具列清楚，并说明怎么用}。不涉及制度评价。

\textbf{一、意定监护协议（最值得优先考虑的一项）}
\begin{itemize}
\item \textbf{法律依据}：《民法典》第三十三条确立了意定监护制度——具有完全民事行为能力的成年人，可以事先与他人协商确定自己的监护人，在本人丧失或部分丧失民事行为能力时，由该监护人履行监护职责；
\item \textbf{对伴侣关系的实际意义}：这份协议可以让伴侣在自己失能、失智或昏迷时，\emph{取得合法的决策与照顾地位}。内容可写入医疗决定、生活照护、财产管理等方面的授权范围；
\item \textbf{怎么办理}：与他方签订书面意定监护协议，建议同时办理\textbf{公证}——公证能显著提高协议在机构（医院、银行等）面前的被认可度与实际可执行性。也可向公证机构咨询具体文本与流程；
\item \textbf{常见误区}：口头约定或仅在微信里说过的"我要是出事你负责"，在实际场景中基本不具有可操作性。书面的、可被第三方识别的文件才有用。
\end{itemize}

\textbf{二、医疗签字与陪诊}
\begin{itemize}
\item 医院在手术同意、紧急救治等环节通常只认可法定亲属或已登记的监护人签字。因此\textbf{意定监护协议是解决签字问题的关键工具}——在协议中明确医疗决策授权，并在就诊前把协议复印件带上；
\item 在协议尚未办理的情况下，仍可通过\textbf{授权委托书}（针对特定事项、特定期间的书面授权）处理部分事务；也可通过病案中登记"紧急联系人"来提升实际联系效率；
\item 现实中更常见的情况是——伴侣被拒于签字台之外。这时最有效的做法是带上任何可证明关系与授权意图的书面材料，并向医院医务部门说明；如情况紧急，医院仍有救急处置优先的原则。
\end{itemize}

\textbf{三、财产与身后安排}
\begin{itemize}
\item \textbf{遗嘱}：通过合法形式的遗嘱（公证遗嘱的证明力最强），可以把个人财产留给伴侣（遗产可通过遗赠方式给到非法定继承人，但对继承顺序与形式要求更严，建议咨询律师或公证员）；
\item \textbf{财产约定}：可以通过书面协议对共同购置的房产、车辆、投资等约定份额与归属——注意，\emph{登记名义与出资事实的证明材料都要保留}（转账记录、出资凭证），这比协议本身更能决定实际结果；
\item \textbf{保险受益人}：人身保险的受益人可由投保人或被保险人指定，这是一条门槛低、见效快的路径；同时可考虑为伴侣配置相应的保障；
\item \textbf{共同债务与账户}：合租、合贷、共同理财的书面记录值得保留——出现纠纷时，这是唯一的凭据。
\end{itemize}

\textbf{四、办理顺序建议}
\begin{itemize}
\item 可以先做\textbf{保险受益人指定}（成本最低、生效最快），再办\textbf{意定监护协议的公证}（关系到医疗与失能照护，价值最高），然后是有实际资产的伴侣办理\textbf{遗嘱与财产约定}；
\item 所有材料建议\textbf{各自保存一份纸质原件}，并让一位可信任的第三方知道存放位置；
\item 若涉及未成年子女、跨地区财产或已有婚姻关系的情况，务必咨询专业律师——每一种情况都有需要另外处理的变量。
\end{itemize}

\noindent\textit{【编者按】}本节只陈述现行法律框架下可用的工具及其办理方式，不对制度本身做任何评价。文中的一般性说明不能替代个案的法律意见；涉及具体办理时，请以公证机构与执业律师的答复为准。
